%% ============================================================================
%% Section: Introduction
%% Included from main.tex via %% ============================================================================
%% Section: Introduction
%% Included from main.tex via %% ============================================================================
%% Section: Introduction
%% Included from main.tex via %% ============================================================================
%% Section: Introduction
%% Included from main.tex via \input{sections/introduction}
%% Editorial notes use \TODO{} / \NOTE{} (defined in main.tex); remove before CR.
%% Numbers marked \TODO are NOT yet measured -- see ICDE2027-outline.md, sec. 13.
%% ============================================================================
\section{Introduction}\label{sec:intro}

Causal reasoning has become central to data-driven decision making across
epidemiology, economics, and the empirical sciences. Analysts increasingly ask
\emph{structural} questions of a causal model: \emph{which variables confound
the effect of a treatment $T$ on an outcome $Y$?}, \emph{does a valid adjustment
set exist, and what is it?}, or \emph{is $X$ independent of $Y$ once we condition
on $Z$?} Each of these is answered not from data alone but from the topology of a
\emph{causal directed acyclic graph} (DAG), whose nodes are variables and whose
edges encode direct causal influence~\cite{pearl2009causality}. Such DAGs are no
longer small, hand-drawn diagrams: they are large, continually revised, and
increasingly stored and versioned inside graph databases alongside the domain
data they describe.

Despite this, causal analysis today is largely divorced from the systems that
store the graph. The prevailing practice is to \emph{export} the DAG into a
standalone library---\textsc{DoWhy}, \textsc{dagitty}, \textsc{pgmpy}, or
\textsc{causal-learn}---compute the answer in a separate process, and discard the
in-memory representation afterwards. This round-trip forfeits precisely the
properties a database is built to provide: durable and concurrent access to a
single authoritative model, declarative querying, integration with co-located
data, and---critically---\emph{incremental} response to a changing graph. When an
edge is added or a variable is reclassified as latent, the entire model must be
re-extracted and every downstream index rebuilt from scratch. No graph DBMS
today treats causal structural queries as a first-class, \emph{indexed} workload.

Closing this gap is not merely a matter of wrapping a library behind a stored
procedure. First, the operators analysts need are not all expressible as plain
graph reachability: \emph{d-separation} and \emph{adjustment-set enumeration}
depend on the conditioning set and on edge orientation, and require
conditioning-aware traversal (the Bayes-Ball procedure) rather than a transitive
closure. Second, answered naively---by traversing the stored graph on every
call---each query costs $O(V{+}E)$, and identification tasks issue many such
queries as they enumerate and test candidate adjustment sets, so per-query
traversal does not scale to large DAGs or high query volume. Third, the index
that accelerates these queries must survive updates: a causal graph under active
curation cannot afford a full recompilation after every edit.

\NOTE{Optional 4th challenge sentence: expressing the operators in a language
that is close to how domain experts phrase causal questions, yet compiles to
physical operators a query engine can execute.}

We present \textbf{Causal4j}, a system that makes causal structural queries a
native, indexed workload inside a property-graph DBMS (Neo4j). Causal4j has
three parts, shown in Fig.~\ref{fig:arch}. (i) A \emph{declarative causal query
language} whose operators---ancestors and descendants, mediators, colliders and
confounders, backdoor paths, d-separation, and adjustment-set
identification---read like the questions analysts already ask, and whose
statements compile to physical operators over the graph store. (ii) A
\emph{compiled causal index}: on compilation the DAG stored as
$(:$\textsf{Variable}$)\text{-}[:$\textsf{CAUSES}$]\!\rightarrow\!(:$\textsf{Variable}$)$
is snapshotted into a dense, integer-indexed layout (CSR adjacency, a
topological order and per-node layer) augmented with precomputed
\emph{ancestor} and \emph{descendant} bitsets; structural queries then reduce
from $O(V{+}E)$ traversals to $O(V/64)$ word operations, and d-separation and
adjustment enumeration are cast as bitset intersections. (iii) \emph{Incremental
maintenance} of that index under edge and variable updates, so a mutating causal
graph is kept queryable without full recompilation.

\NOTE{Insert Fig.~1 (architecture) here -- provided by figures/architecture.tex.}

This paper makes the following contributions:
\begin{itemize}
  \item A \textbf{declarative causal query language}, embedded in a graph DBMS,
  covering structural queries, identification (adjustment sets, backdoor
  criterion), and a quantitative layer (interventions, counterfactuals), with a
  defined compilation from language operators to physical graph operators
  (\S\ref{sec:lang}).
  \item A \textbf{compiled reachability index for causal DAGs}---dense CSR
  adjacency with topological-layer pruning and precomputed ancestor/descendant
  bitsets---that reduces structural queries from $O(V{+}E)$ to $O(V/64)$ and
  supports bitset Bayes-Ball d-separation and adjustment-set enumeration
  (\S\ref{sec:index}).
  \item An \textbf{incremental maintenance} scheme that updates the index under
  edge and variable mutations, bounding work to the affected region rather than
  recompiling the graph (\S\ref{sec:incr}).
  \item An \textbf{empirical study} on standard Bayesian-network DAGs and
  synthetic scale-free DAGs, comparing Causal4j against in-DBMS Cypher traversal,
  Neo4j GDS, and external causal libraries on latency, throughput, index build
  time, memory, and update cost, released as a reproducible artifact
  (\S\ref{sec:eval}).
\end{itemize}

\NOTE{Results preview -- fill once the eval in \S\ref{sec:eval} exists. Keep it
one or two sentences with the headline numbers.}
On DAGs ranging from \TODO{$V_{\min}$} to \TODO{$V_{\max}$} variables, Causal4j
answers d-separation and adjustment-set queries \TODO{$N\times$} faster than
in-DBMS Cypher traversal and \TODO{$M\times$} faster than round-tripping to an
external library, builds its index in \TODO{under X\,ms}, and applies an edge
update in \TODO{Y}, versus \TODO{Z} for a full recompilation.

\emph{Scope note for reviewers.} This is a data-engineering contribution: a
declarative query language embedded in a graph DBMS, a physical reachability index
for causal DAGs, query processing over that index, and its incremental
maintenance under updates. The structural-causal-model semantics (d-separation,
the backdoor criterion, adjustment sets) are the \emph{workload} the engine
supports, drawn from established causal theory~\cite{pearl2009causality}; they are
not claimed as a contribution. What is new is treating that workload as a
first-class, indexed, incrementally maintained query workload \emph{inside} the
store, rather than in an external library.

The remainder of the paper is organized as follows. \S\ref{sec:bg} defines
causal DAGs, the query catalog, and the in-DBMS problem. \S\ref{sec:lang}
presents the language and its compilation; \S\ref{sec:compile} the compilation
of the causal snapshot; \S\ref{sec:index} the index and query processing;
\S\ref{sec:incr} incremental maintenance; and \S\ref{sec:quant} the quantitative
layer. \S\ref{sec:eval} reports the evaluation, \S\ref{sec:related} surveys
related work, and \S\ref{sec:concl} concludes.

%% Editorial notes use \TODO{} / \NOTE{} (defined in main.tex); remove before CR.
%% Numbers marked \TODO are NOT yet measured -- see ICDE2027-outline.md, sec. 13.
%% ============================================================================
\section{Introduction}\label{sec:intro}

Causal reasoning has become central to data-driven decision making across
epidemiology, economics, and the empirical sciences. Analysts increasingly ask
\emph{structural} questions of a causal model: \emph{which variables confound
the effect of a treatment $T$ on an outcome $Y$?}, \emph{does a valid adjustment
set exist, and what is it?}, or \emph{is $X$ independent of $Y$ once we condition
on $Z$?} Each of these is answered not from data alone but from the topology of a
\emph{causal directed acyclic graph} (DAG), whose nodes are variables and whose
edges encode direct causal influence~\cite{pearl2009causality}. Such DAGs are no
longer small, hand-drawn diagrams: they are large, continually revised, and
increasingly stored and versioned inside graph databases alongside the domain
data they describe.

Despite this, causal analysis today is largely divorced from the systems that
store the graph. The prevailing practice is to \emph{export} the DAG into a
standalone library---\textsc{DoWhy}, \textsc{dagitty}, \textsc{pgmpy}, or
\textsc{causal-learn}---compute the answer in a separate process, and discard the
in-memory representation afterwards. This round-trip forfeits precisely the
properties a database is built to provide: durable and concurrent access to a
single authoritative model, declarative querying, integration with co-located
data, and---critically---\emph{incremental} response to a changing graph. When an
edge is added or a variable is reclassified as latent, the entire model must be
re-extracted and every downstream index rebuilt from scratch. No graph DBMS
today treats causal structural queries as a first-class, \emph{indexed} workload.

Closing this gap is not merely a matter of wrapping a library behind a stored
procedure. First, the operators analysts need are not all expressible as plain
graph reachability: \emph{d-separation} and \emph{adjustment-set enumeration}
depend on the conditioning set and on edge orientation, and require
conditioning-aware traversal (the Bayes-Ball procedure) rather than a transitive
closure. Second, answered naively---by traversing the stored graph on every
call---each query costs $O(V{+}E)$, and identification tasks issue many such
queries as they enumerate and test candidate adjustment sets, so per-query
traversal does not scale to large DAGs or high query volume. Third, the index
that accelerates these queries must survive updates: a causal graph under active
curation cannot afford a full recompilation after every edit.

\NOTE{Optional 4th challenge sentence: expressing the operators in a language
that is close to how domain experts phrase causal questions, yet compiles to
physical operators a query engine can execute.}

We present \textbf{Causal4j}, a system that makes causal structural queries a
native, indexed workload inside a property-graph DBMS (Neo4j). Causal4j has
three parts, shown in Fig.~\ref{fig:arch}. (i) A \emph{declarative causal query
language} whose operators---ancestors and descendants, mediators, colliders and
confounders, backdoor paths, d-separation, and adjustment-set
identification---read like the questions analysts already ask, and whose
statements compile to physical operators over the graph store. (ii) A
\emph{compiled causal index}: on compilation the DAG stored as
$(:$\textsf{Variable}$)\text{-}[:$\textsf{CAUSES}$]\!\rightarrow\!(:$\textsf{Variable}$)$
is snapshotted into a dense, integer-indexed layout (CSR adjacency, a
topological order and per-node layer) augmented with precomputed
\emph{ancestor} and \emph{descendant} bitsets; structural queries then reduce
from $O(V{+}E)$ traversals to $O(V/64)$ word operations, and d-separation and
adjustment enumeration are cast as bitset intersections. (iii) \emph{Incremental
maintenance} of that index under edge and variable updates, so a mutating causal
graph is kept queryable without full recompilation.

\NOTE{Insert Fig.~1 (architecture) here -- provided by figures/architecture.tex.}

This paper makes the following contributions:
\begin{itemize}
  \item A \textbf{declarative causal query language}, embedded in a graph DBMS,
  covering structural queries, identification (adjustment sets, backdoor
  criterion), and a quantitative layer (interventions, counterfactuals), with a
  defined compilation from language operators to physical graph operators
  (\S\ref{sec:lang}).
  \item A \textbf{compiled reachability index for causal DAGs}---dense CSR
  adjacency with topological-layer pruning and precomputed ancestor/descendant
  bitsets---that reduces structural queries from $O(V{+}E)$ to $O(V/64)$ and
  supports bitset Bayes-Ball d-separation and adjustment-set enumeration
  (\S\ref{sec:index}).
  \item An \textbf{incremental maintenance} scheme that updates the index under
  edge and variable mutations, bounding work to the affected region rather than
  recompiling the graph (\S\ref{sec:incr}).
  \item An \textbf{empirical study} on standard Bayesian-network DAGs and
  synthetic scale-free DAGs, comparing Causal4j against in-DBMS Cypher traversal,
  Neo4j GDS, and external causal libraries on latency, throughput, index build
  time, memory, and update cost, released as a reproducible artifact
  (\S\ref{sec:eval}).
\end{itemize}

\NOTE{Results preview -- fill once the eval in \S\ref{sec:eval} exists. Keep it
one or two sentences with the headline numbers.}
On DAGs ranging from \TODO{$V_{\min}$} to \TODO{$V_{\max}$} variables, Causal4j
answers d-separation and adjustment-set queries \TODO{$N\times$} faster than
in-DBMS Cypher traversal and \TODO{$M\times$} faster than round-tripping to an
external library, builds its index in \TODO{under X\,ms}, and applies an edge
update in \TODO{Y}, versus \TODO{Z} for a full recompilation.

\emph{Scope note for reviewers.} This is a data-engineering contribution: a
declarative query language embedded in a graph DBMS, a physical reachability index
for causal DAGs, query processing over that index, and its incremental
maintenance under updates. The structural-causal-model semantics (d-separation,
the backdoor criterion, adjustment sets) are the \emph{workload} the engine
supports, drawn from established causal theory~\cite{pearl2009causality}; they are
not claimed as a contribution. What is new is treating that workload as a
first-class, indexed, incrementally maintained query workload \emph{inside} the
store, rather than in an external library.

The remainder of the paper is organized as follows. \S\ref{sec:bg} defines
causal DAGs, the query catalog, and the in-DBMS problem. \S\ref{sec:lang}
presents the language and its compilation; \S\ref{sec:compile} the compilation
of the causal snapshot; \S\ref{sec:index} the index and query processing;
\S\ref{sec:incr} incremental maintenance; and \S\ref{sec:quant} the quantitative
layer. \S\ref{sec:eval} reports the evaluation, \S\ref{sec:related} surveys
related work, and \S\ref{sec:concl} concludes.

%% Editorial notes use \TODO{} / \NOTE{} (defined in main.tex); remove before CR.
%% Numbers marked \TODO are NOT yet measured -- see ICDE2027-outline.md, sec. 13.
%% ============================================================================
\section{Introduction}\label{sec:intro}

Causal reasoning has become central to data-driven decision making across
epidemiology, economics, and the empirical sciences. Analysts increasingly ask
\emph{structural} questions of a causal model: \emph{which variables confound
the effect of a treatment $T$ on an outcome $Y$?}, \emph{does a valid adjustment
set exist, and what is it?}, or \emph{is $X$ independent of $Y$ once we condition
on $Z$?} Each of these is answered not from data alone but from the topology of a
\emph{causal directed acyclic graph} (DAG), whose nodes are variables and whose
edges encode direct causal influence~\cite{pearl2009causality}. Such DAGs are no
longer small, hand-drawn diagrams: they are large, continually revised, and
increasingly stored and versioned inside graph databases alongside the domain
data they describe.

Despite this, causal analysis today is largely divorced from the systems that
store the graph. The prevailing practice is to \emph{export} the DAG into a
standalone library---\textsc{DoWhy}, \textsc{dagitty}, \textsc{pgmpy}, or
\textsc{causal-learn}---compute the answer in a separate process, and discard the
in-memory representation afterwards. This round-trip forfeits precisely the
properties a database is built to provide: durable and concurrent access to a
single authoritative model, declarative querying, integration with co-located
data, and---critically---\emph{incremental} response to a changing graph. When an
edge is added or a variable is reclassified as latent, the entire model must be
re-extracted and every downstream index rebuilt from scratch. No graph DBMS
today treats causal structural queries as a first-class, \emph{indexed} workload.

Closing this gap is not merely a matter of wrapping a library behind a stored
procedure. First, the operators analysts need are not all expressible as plain
graph reachability: \emph{d-separation} and \emph{adjustment-set enumeration}
depend on the conditioning set and on edge orientation, and require
conditioning-aware traversal (the Bayes-Ball procedure) rather than a transitive
closure. Second, answered naively---by traversing the stored graph on every
call---each query costs $O(V{+}E)$, and identification tasks issue many such
queries as they enumerate and test candidate adjustment sets, so per-query
traversal does not scale to large DAGs or high query volume. Third, the index
that accelerates these queries must survive updates: a causal graph under active
curation cannot afford a full recompilation after every edit.

\NOTE{Optional 4th challenge sentence: expressing the operators in a language
that is close to how domain experts phrase causal questions, yet compiles to
physical operators a query engine can execute.}

We present \textbf{Causal4j}, a system that makes causal structural queries a
native, indexed workload inside a property-graph DBMS (Neo4j). Causal4j has
three parts, shown in Fig.~\ref{fig:arch}. (i) A \emph{declarative causal query
language} whose operators---ancestors and descendants, mediators, colliders and
confounders, backdoor paths, d-separation, and adjustment-set
identification---read like the questions analysts already ask, and whose
statements compile to physical operators over the graph store. (ii) A
\emph{compiled causal index}: on compilation the DAG stored as
$(:$\textsf{Variable}$)\text{-}[:$\textsf{CAUSES}$]\!\rightarrow\!(:$\textsf{Variable}$)$
is snapshotted into a dense, integer-indexed layout (CSR adjacency, a
topological order and per-node layer) augmented with precomputed
\emph{ancestor} and \emph{descendant} bitsets; structural queries then reduce
from $O(V{+}E)$ traversals to $O(V/64)$ word operations, and d-separation and
adjustment enumeration are cast as bitset intersections. (iii) \emph{Incremental
maintenance} of that index under edge and variable updates, so a mutating causal
graph is kept queryable without full recompilation.

\NOTE{Insert Fig.~1 (architecture) here -- provided by figures/architecture.tex.}

This paper makes the following contributions:
\begin{itemize}
  \item A \textbf{declarative causal query language}, embedded in a graph DBMS,
  covering structural queries, identification (adjustment sets, backdoor
  criterion), and a quantitative layer (interventions, counterfactuals), with a
  defined compilation from language operators to physical graph operators
  (\S\ref{sec:lang}).
  \item A \textbf{compiled reachability index for causal DAGs}---dense CSR
  adjacency with topological-layer pruning and precomputed ancestor/descendant
  bitsets---that reduces structural queries from $O(V{+}E)$ to $O(V/64)$ and
  supports bitset Bayes-Ball d-separation and adjustment-set enumeration
  (\S\ref{sec:index}).
  \item An \textbf{incremental maintenance} scheme that updates the index under
  edge and variable mutations, bounding work to the affected region rather than
  recompiling the graph (\S\ref{sec:incr}).
  \item An \textbf{empirical study} on standard Bayesian-network DAGs and
  synthetic scale-free DAGs, comparing Causal4j against in-DBMS Cypher traversal,
  Neo4j GDS, and external causal libraries on latency, throughput, index build
  time, memory, and update cost, released as a reproducible artifact
  (\S\ref{sec:eval}).
\end{itemize}

\NOTE{Results preview -- fill once the eval in \S\ref{sec:eval} exists. Keep it
one or two sentences with the headline numbers.}
On DAGs ranging from \TODO{$V_{\min}$} to \TODO{$V_{\max}$} variables, Causal4j
answers d-separation and adjustment-set queries \TODO{$N\times$} faster than
in-DBMS Cypher traversal and \TODO{$M\times$} faster than round-tripping to an
external library, builds its index in \TODO{under X\,ms}, and applies an edge
update in \TODO{Y}, versus \TODO{Z} for a full recompilation.

\emph{Scope note for reviewers.} This is a data-engineering contribution: a
declarative query language embedded in a graph DBMS, a physical reachability index
for causal DAGs, query processing over that index, and its incremental
maintenance under updates. The structural-causal-model semantics (d-separation,
the backdoor criterion, adjustment sets) are the \emph{workload} the engine
supports, drawn from established causal theory~\cite{pearl2009causality}; they are
not claimed as a contribution. What is new is treating that workload as a
first-class, indexed, incrementally maintained query workload \emph{inside} the
store, rather than in an external library.

The remainder of the paper is organized as follows. \S\ref{sec:bg} defines
causal DAGs, the query catalog, and the in-DBMS problem. \S\ref{sec:lang}
presents the language and its compilation; \S\ref{sec:compile} the compilation
of the causal snapshot; \S\ref{sec:index} the index and query processing;
\S\ref{sec:incr} incremental maintenance; and \S\ref{sec:quant} the quantitative
layer. \S\ref{sec:eval} reports the evaluation, \S\ref{sec:related} surveys
related work, and \S\ref{sec:concl} concludes.

%% Editorial notes use \TODO{} / \NOTE{} (defined in main.tex); remove before CR.
%% Numbers marked \TODO are NOT yet measured -- see ICDE2027-outline.md, sec. 13.
%% ============================================================================
\section{Introduction}\label{sec:intro}

Causal reasoning has become central to data-driven decision making across
epidemiology, economics, and the empirical sciences. Analysts increasingly ask
\emph{structural} questions of a causal model: \emph{which variables confound
the effect of a treatment $T$ on an outcome $Y$?}, \emph{does a valid adjustment
set exist, and what is it?}, or \emph{is $X$ independent of $Y$ once we condition
on $Z$?} Each of these is answered not from data alone but from the topology of a
\emph{causal directed acyclic graph} (DAG), whose nodes are variables and whose
edges encode direct causal influence~\cite{pearl2009causality}. Such DAGs are no
longer small, hand-drawn diagrams: they are large, continually revised, and
increasingly stored and versioned inside graph databases alongside the domain
data they describe.

Despite this, causal analysis today is largely divorced from the systems that
store the graph. The prevailing practice is to \emph{export} the DAG into a
standalone library---\textsc{DoWhy}, \textsc{dagitty}, \textsc{pgmpy}, or
\textsc{causal-learn}---compute the answer in a separate process, and discard the
in-memory representation afterwards. This round-trip forfeits precisely the
properties a database is built to provide: durable and concurrent access to a
single authoritative model, declarative querying, integration with co-located
data, and---critically---\emph{incremental} response to a changing graph. When an
edge is added or a variable is reclassified as latent, the entire model must be
re-extracted and every downstream index rebuilt from scratch. No graph DBMS
today treats causal structural queries as a first-class, \emph{indexed} workload.

Closing this gap is not merely a matter of wrapping a library behind a stored
procedure. First, the operators analysts need are not all expressible as plain
graph reachability: \emph{d-separation} and \emph{adjustment-set enumeration}
depend on the conditioning set and on edge orientation, and require
conditioning-aware traversal (the Bayes-Ball procedure) rather than a transitive
closure. Second, answered naively---by traversing the stored graph on every
call---each query costs $O(V{+}E)$, and identification tasks issue many such
queries as they enumerate and test candidate adjustment sets, so per-query
traversal does not scale to large DAGs or high query volume. Third, the index
that accelerates these queries must survive updates: a causal graph under active
curation cannot afford a full recompilation after every edit.

\NOTE{Optional 4th challenge sentence: expressing the operators in a language
that is close to how domain experts phrase causal questions, yet compiles to
physical operators a query engine can execute.}

We present \textbf{Causal4j}, a system that makes causal structural queries a
native, indexed workload inside a property-graph DBMS (Neo4j). Causal4j has
three parts, shown in Fig.~\ref{fig:arch}. (i) A \emph{declarative causal query
language} whose operators---ancestors and descendants, mediators, colliders and
confounders, backdoor paths, d-separation, and adjustment-set
identification---read like the questions analysts already ask, and whose
statements compile to physical operators over the graph store. (ii) A
\emph{compiled causal index}: on compilation the DAG stored as
$(:$\textsf{Variable}$)\text{-}[:$\textsf{CAUSES}$]\!\rightarrow\!(:$\textsf{Variable}$)$
is snapshotted into a dense, integer-indexed layout (CSR adjacency, a
topological order and per-node layer) augmented with precomputed
\emph{ancestor} and \emph{descendant} bitsets; structural queries then reduce
from $O(V{+}E)$ traversals to $O(V/64)$ word operations, and d-separation and
adjustment enumeration are cast as bitset intersections. (iii) \emph{Incremental
maintenance} of that index under edge and variable updates, so a mutating causal
graph is kept queryable without full recompilation.

\NOTE{Insert Fig.~1 (architecture) here -- provided by figures/architecture.tex.}

This paper makes the following contributions:
\begin{itemize}
  \item A \textbf{declarative causal query language}, embedded in a graph DBMS,
  covering structural queries, identification (adjustment sets, backdoor
  criterion), and a quantitative layer (interventions, counterfactuals), with a
  defined compilation from language operators to physical graph operators
  (\S\ref{sec:lang}).
  \item A \textbf{compiled reachability index for causal DAGs}---dense CSR
  adjacency with topological-layer pruning and precomputed ancestor/descendant
  bitsets---that reduces structural queries from $O(V{+}E)$ to $O(V/64)$ and
  supports bitset Bayes-Ball d-separation and adjustment-set enumeration
  (\S\ref{sec:index}).
  \item An \textbf{incremental maintenance} scheme that updates the index under
  edge and variable mutations, bounding work to the affected region rather than
  recompiling the graph (\S\ref{sec:incr}).
  \item An \textbf{empirical study} on standard Bayesian-network DAGs and
  synthetic scale-free DAGs, comparing Causal4j against in-DBMS Cypher traversal,
  Neo4j GDS, and external causal libraries on latency, throughput, index build
  time, memory, and update cost, released as a reproducible artifact
  (\S\ref{sec:eval}).
\end{itemize}

\NOTE{Results preview -- fill once the eval in \S\ref{sec:eval} exists. Keep it
one or two sentences with the headline numbers.}
On DAGs ranging from \TODO{$V_{\min}$} to \TODO{$V_{\max}$} variables, Causal4j
answers d-separation and adjustment-set queries \TODO{$N\times$} faster than
in-DBMS Cypher traversal and \TODO{$M\times$} faster than round-tripping to an
external library, builds its index in \TODO{under X\,ms}, and applies an edge
update in \TODO{Y}, versus \TODO{Z} for a full recompilation.

\emph{Scope note for reviewers.} This is a data-engineering contribution: a
declarative query language embedded in a graph DBMS, a physical reachability index
for causal DAGs, query processing over that index, and its incremental
maintenance under updates. The structural-causal-model semantics (d-separation,
the backdoor criterion, adjustment sets) are the \emph{workload} the engine
supports, drawn from established causal theory~\cite{pearl2009causality}; they are
not claimed as a contribution. What is new is treating that workload as a
first-class, indexed, incrementally maintained query workload \emph{inside} the
store, rather than in an external library.

The remainder of the paper is organized as follows. \S\ref{sec:bg} defines
causal DAGs, the query catalog, and the in-DBMS problem. \S\ref{sec:lang}
presents the language and its compilation; \S\ref{sec:compile} the compilation
of the causal snapshot; \S\ref{sec:index} the index and query processing;
\S\ref{sec:incr} incremental maintenance; and \S\ref{sec:quant} the quantitative
layer. \S\ref{sec:eval} reports the evaluation, \S\ref{sec:related} surveys
related work, and \S\ref{sec:concl} concludes.
